\centerline{\bf \Huge Геометрия II.}
\centerline{\bf \Huge Вписанные углы II}

\begin{flushright}
        \scriptsize{}
\end{flushright}

\problem{1}\textbf{Вспомогательная для 4 задачи из прошлого листочка.}
На окружности взяты точки $A, C_1, B, A_1, C, B_1$ в указанном
порядке.

а) Докажите, что если прямые $AA1, BB1$ и $CC1$ являются биссектрисами углов треугольника $ABC$, то они являются высотами треугольника $A_1B_1C_1.$

б) Докажите, что если прямые $AA1, BB1$ и $CC1$ являются высотами
треугольника $ABC$, то они являются биссектрисами углов треугольника $A_1B_1C_1.$

\problem{2}
Окружность $S_1$ с диаметром AB пересекает окружность $S_2$
с центром $A$ в точках $C$ и $D$. Через точку B проведена прямая,
пересекающая $S_2$ в точке $M$, лежащей внутри $S_1$, а $S_1$ в точке $N$.
Докажите, что $MN^2 = CN \cdot ND.$

\problem{3}
Продолжение биссектрисы $AD$ остроугольного треугольника $ABC$ пересекает описанную окружность в точке $E$. Из точки $D$ на стороны $AB$ и $AC$ опущены перпендикуляры $DP$ и $DQ$. Докажите, что $S_{ABC} =S_{APEQ}$.

\problem{4}
Четыре прямые образуют четыре треугольника.

а) Докажите, что описанные окружности этих треугольников имеют общую точку \textit{(точка Микеля)}.

б) Докажите, что центры описанных окружностей этих треугольников лежат на одной окружности, проходящей через точку Микеля.

\problem{5}\textbf{Затравка к следующему (тематическому) листочку}
В треугольнике $ABC$ проведены медианы $AF$ и $CE$. Докажите, что если $\angle BAF=\angle BCE = 30^{\circ}$, то треугольник $ABC$ правильный.